\documentclass[12pt]{article}
\usepackage[utf8]{inputenc}
\usepackage[T1]{fontenc}
\usepackage[margin=1in]{geometry}
\usepackage{setspace}
\usepackage{xcolor}
\usepackage[colorlinks=true, urlcolor=blue!60!black, linkcolor=blue!60!black]{hyperref}
\setlength{\parindent}{2em}
\setlength{\parskip}{1.25em}
\renewcommand{\baselinestretch}{1.0}
\begin{document}
\begin{center}
{\LARGE \textit{United States v. E.C. Knight Co.}, 156 U.S. 1 (1895)\par}
\vspace{0.8em}
{\large Case Brief --- Four-Page Edition\par}
\end{center}
\vspace{1em}

\section{Memory Jogger}

The Sugar Trust's purchase of four Philadelphia refineries gave it control of about 98\% of American sugar refining, yet the Court held 8--1 that the Sherman Act did not reach the acquisition, because the combination controlled local manufacture and, on the facts alleged, affected interstate commerce only indirectly. Justice Harlan dissented alone, and the manufacture/commerce line was later repudiated (\href{https://supreme.justia.com/cases/federal/us/156/1/}{156 U.S. 1}).

\section{Detailed Case Facts}

In March 1892, the American Sugar Refining Company, a New Jersey corporation that already controlled most American sugar refining, contracted to buy the stock of four competing Philadelphia refineries: the E.C. Knight Company, the Franklin Sugar Refining Company, Spreckels' Sugar Refining Company, and the Delaware Sugar House.Only the Revere Sugar Refinery of Boston, with about 2\% of output, remained independent. The ``Sugar Trust'' thus controlled approximately 98\% of all refined sugar manufactured in the United States, a household necessity sold nationwide.

\section{Procedural History}

The United States filed a bill in equity in the Circuit Court for the Eastern District of Pennsylvania under the Sherman Antitrust Act of 1890, alleging that the stock purchases were a contract, combination, or conspiracy to restrain and monopolize interstate and foreign trade in refined sugar, and asking the court to cancel the agreements. The Circuit Court dismissed the bill, concluding that the contracts concerned only manufacturing within Pennsylvania, not interstate commerce. The Third Circuit affirmed, and the United States appealed.

\section{Judicial Votes}

Affirmed, \textbf{8--1}.
\begin{itemize}
\item \textbf{Majority:} Chief Justice Fuller, joined by Justices Field, Gray, Brewer, Brown, Shiras, Jackson, and White.
\item \textbf{Concurring:} None.
\item \textbf{Dissenting:} Justice Harlan, alone.
\end{itemize}

\section{Holding}

The Sherman Act, as an exercise of the commerce power, could not reach a monopoly in the manufacture of refined sugar, even a near-complete one, where the acquired refineries were located in a single state. Manufacturing precedes interstate distribution and affects interstate commerce only incidentally and indirectly, so regulating manufacturing monopolies belonged to the states' police power, not to Congress. On the government's proofs, which showed an acquisition of manufacturing stock rather than a direct restraint on interstate sales or transportation, the bill was properly dismissed.

\section{Analysis of Opinions}

\textbf{Majority (Fuller).} Fuller rested the Court's dual-federalism holding on four formalist distinctions.

\emph{Manufacture versus commerce.} Drawing on \textit{Kidd v.\ Pearson}, Fuller explained that ``manufacture is transformation---the fashioning of raw materials into a change of form for use,'' whereas commerce consists of buying, selling, and interstate transportation. Even if every pound of sugar refined in Philadelphia was destined for shipment to other states, manufacture came first and stood apart from the commerce that followed.

\emph{Exclusive spheres.} Treating state and federal authority as mutually exclusive, Fuller assigned whatever belongs to commerce to the United States and whatever does not to the police power of the states. 

\emph{Direct versus indirect effects.} Fuller conceded that controlling 98\% of manufacture involves, in a certain sense, control of the product's disposition, but called that a secondary sense. The combination affected interstate commerce only incidentally and indirectly, and the government had shown an acquisition of manufacturing stock, not a direct restraint on interstate trade.

\emph{The slippery slope.} If Congress could regulate local manufacture whenever its products entered interstate trade, Fuller warned, ``comparatively little of business operations and affairs would be left for state control,'' and agriculture and mining would fall under federal supervision.

\textbf{Harlan (dissenting).} In a prophetic solo dissent, Harlan rejected the rigid boundary between production and trade. Relying on \textit{Gibbons v.\ Ogden} and \textit{McCulloch v.\ Maryland}, he argued that interstate commerce is not mere transportation but includes the purchase and sale of goods intended for shipment among the states. A trust controlling 98\% of a necessity refined for nationwide sale directly obstructed interstate commerce and raised prices for consumers in every state. Because no single state could police a multistate combination chartered in New Jersey, Congress had power to act.

\textbf{Later developments.} Holmes's current-of-commerce doctrine in \textit{Swift} (1905) distinguished \textit{E.C. Knight}; \textit{Jones \& Laughlin} (1937) and \textit{Wickard} (1942) buried the manufacture/commerce wall.

\section{Examples: Future Applications}

\begin{enumerate}
\item \textbf{Same-Side (Beyond the Power): A Homemade Tool Near a Park.}  Congress bans purely local, noncommercial possession of a homemade tool near a park. Under \textit{Lopez} and \textit{Morrison}, noneconomic conduct cannot be aggregated into a general police power.
\item \textbf{Same-Side (Beyond the Power): A Bicycle Mandate.} Congress requires every adult to buy an American-made electric bicycle.\textit{NFIB v.\ Sebelius} supplies a modern limit: Congress may regulate existing commercial activity, not compel people to enter commerce.
\item \textbf{Opposite-Side (Within the Power): A Chip-Fabrication Monopoly.} A holding company buys four Ohio chip plants, gaining 95\% of a chip sold nationwide.Under \textit{Jones \& Laughlin}, \textit{Wickard}, and \textit{Raich}, Congress may reach the merger.
\item \textbf{Opposite-Side (Within the Power): Grain-Elevator Price-Fixing.} Kansas elevators fix the prices they pay for wheat shipped to other states; \textit{Swift} reached similar stockyard price-fixing.
\item \textbf{Fence-Sitter: Homemade Firearms.} A ban on possessing unserialized homemade guns is part of a broad scheme regulating the interstate market in parts and kits. Possession is noneconomic, as in \textit{Lopez}, but the market scheme resembles \textit{Gonzales v.\ Raich}.
\end{enumerate}

\section{Critique}

\textbf{Formalism divorced from economics.} \textit{E.C. Knight} is a textbook study in the pitfalls of conceptual formalism divorced from economic reality. As Harlan showed, monopolizing 98\% of a nationally consumed commodity's supply is functionally indistinguishable from restraining its interstate sale and price.

\textbf{The search for limits.} Yet students should also ask why the Court did not abandon all limits after 1937. \textit{E.C. Knight}'s failure forced later Courts, in \textit{Lopez}, \textit{Morrison}, and \textit{NFIB}, to search for a more workable boundary, such as the economic/noneconomic and activity/inactivity distinctions, rather than resurrect the unworkable manufacture/commerce divide.

\section{Key Quotations}

\begin{quote}
``Commerce succeeds to manufacture, and is not a part of it. \ldots That which belongs to commerce is within the jurisdiction of the United States, but that which does not belong to commerce is within the jurisdiction of the police power of the state.''
\end{quote}
\noindent --- Chief Justice Fuller, Opinion of the Court, 156 U.S. at 12.

\begin{quote}
``Doubtless the power to control the manufacture of a given thing involves, in a certain sense, the control of its disposition, but this is a secondary, and not the primary, sense, and although the exercise of that power may result in bringing the operation of commerce into play, it does not control it, and affects it only incidentally and indirectly.''
\end{quote}
\noindent --- Chief Justice Fuller, Opinion of the Court, 156 U.S. at 12.

\begin{quote}
``Nevertheless it does not follow that an attempt to monopolize, or the actual monopoly of, the manufacture was an attempt, whether executory or consummated, to monopolize commerce, even though, in order to dispose of the product, the instrumentality of commerce was necessarily invoked.''
\end{quote}
\noindent --- Chief Justice Fuller, Opinion of the Court, 156 U.S. at 17.

\end{document}
